\begin{definition}[Classes laterais e índice]
  Se $H\leq G$, definimos uma relação em $G$ da seguinte maneira
  \begin{align*}
    \text{Se }a,b\in G,a\equiv_H b \text{ sé y so sé }ab^{-1}\in H.
  \end{align*}
  Temos
  \begin{itemize}
    \item Esta relação é uma relação de equivalência en $G$.
    \item A classe de equivalência de um elemento $a\in G$ é dada por
      \begin{align*}
        \{b\in G:B\equiv_H a\}&=\{b\in G: ba^{-1}=h\text{ para algúm }h\in H\}\\
        &=\{b\in G: b=ha\text{ para algúm }h\in H\}\\
        &=\{ha:h\in H\}:=Ha
      \end{align*}
      e denominada uma \textbf{classe lateral} (á dereita) de $H$ em $G$ (que contem $a$).
    \item Se $a,b\in G$, então $Ha=Hb$ ou $Ha\cap  Hb=\emptyset$.
    \item A cardinalidade do conjunto das classes laterais (á dereita) de $H$ em $G$ é denominada o \textbf{índice} de $H$ em $G$ e denotada por $[G:H]$.
    \item Para todo $a\in G$ temos $|Ha| = |H|$. 
  \end{itemize}
\end{definition}
\begin{theorem}[Teorema de Lagrange]
  Se $G$ é um grupo finito e $H\leq G$. Então 
  $$|G|=[G:H]|H|,$$
  em particular $|H|$ e $[G:H]$ dividem $|G|$. 
\end{theorem}
\begin{corollary}
  Se $|G|=p$, onde $p$ é um primo, então $G$ é um grupo cíclico. 
\end{corollary}
\begin{corollary}
  Sejam $G$ um grupo finito e $a\in G$. Então $o(a)$ divide $|G|$. 
\end{corollary}
\begin{corollary}
  Se $H,K\leq G$ e $\mdc(|H|,|K|)=1$ então $H\cap K=\{1\}$. 
\end{corollary}
Poderíamos também definir a seguinte relação de equivalência em $G$
\begin{align*}
  \text{Se }a,b\in G, a\equiv^{H} b\text{ sé e só sé }a^{-1}b\in H.
\end{align*}
Neste caso, utilizando as mesmas idéias anteriores, temos um conjunto de classes laterais (\textbf{à esquerda}) de $H$ em $G$, $aH$ onde $a\in G$, com propriedades análogas às da classe lateral à dereita.
\begin{note}\label{obs:propclasseslaterais}
  Note que se $G$ é um grupo e $H\leq G$
  \begin{itemize}
    \item Se $G$ é abeliano então $aH=Ha$, para todo $a\in G$.
    \item Se $[G:H]=2$ então $aH=Ha$, para todo $a\in G$.
    \item No entanto, em geral, podemos ter $aH\neq Ha$, para algúm $a\in G$. 
  \end{itemize}
\end{note}
Tome $G=GL_2(\mathbb{R})$, logo tome $H=\langle a \rangle$ en donde $a$ é
\begin{align*}
  a:=\begin{pmatrix}
    -1 & 0 \\
    0 & 1
  \end{pmatrix}.
\end{align*}
Note que $o(a)=2$, logo $\langle a \rangle=\{1,A\}$.

Agora seja $b\in GL_2(\mathbb{R})$ da seguinte maneira
\begin{align*}
  b:=\begin{pmatrix}
    1 & 1 \\
    0 & 1
  \end{pmatrix}
\end{align*}
Logo
\begin{align*}
  bH&=\{b,ba\}\\
  Hb&=\{b,ab\}.
\end{align*}
Assim, se você realizar os seguintes calculos, vera que
\begin{align*}
  ab&=\begin{pmatrix}
    -1 & -1 \\
    0 & 1
  \end{pmatrix}\\
  ba&=\begin{pmatrix}
    -1 & 1 \\
    0 & 1
  \end{pmatrix}.
\end{align*}
Portanto $bH\neq Hb$ con $b\in Gl_2(\mathbb{R})$ como mencionamos anteriormente.
